\chapter*{Lời mở đầu}
\pagestyle{plain}
\enlargethispage{3\baselineskip}
\addcontentsline{toc}{chapter}{Lời mở đầu}

\section*{1. Tính cấp thiết của Đề tài}
Trong các trung tâm dữ liệu (Data Center) và phòng máy chủ, việc quản lý và giám sát thiết bị phần cứng đóng vai trò rất quan trọng. Tuy nhiên, các dữ liệu giám sát (CPU/RAM, trạng thái máy chủ, lịch sử bảo trì) thường bị tách rời khỏi vị trí vật lý thực tế. Khi xảy ra sự cố, kỹ thuật viên gặp khó khăn trong việc xác định chính xác máy chủ bị lỗi giữa các tủ Rack, gây tốn thời gian xử lý.

Sự kết hợp giữa công nghệ \textbf{Thực tế Tăng cường (AR)} và \textbf{Cơ sở dữ liệu tập trung} cho phép gắn kết dữ liệu vận hành trực tiếp lên thiết bị thực tế qua mã QR/ArUco Marker. Để hệ thống AR và Web Admin truy xuất dữ liệu nhất quán, việc xây dựng một Cơ sở Dữ liệu (CSDL) chuẩn hóa là bước nền tảng cốt lõi.

Nhận thức được điều đó, nhóm chúng em đã lựa chọn đề tài: \textbf{“Thiết kế Cơ sở Dữ liệu cho Hệ thống Giám sát và Bảo trì Cơ sở Hạ tầng Tích hợp Thực tế Tăng cường (AR-IMMS)”} làm nội dung nghiên cứu cho bài tập lớn môn học.

\section*{2. Mục tiêu và Nội dung Thực hiện}
Đề tài tập trung áp dụng các kiến thức môn Thiết kế Cơ sở Dữ liệu để giải quyết bài toán quản lý dữ liệu cho hệ thống AR-IMMS với các mục tiêu chính:
\begin{itemize}
    \item \textbf{Phân tích \& Mô hình hóa Nghiệp vụ:} Khảo sát dòng chảy dữ liệu quản lý hạ tầng phòng máy (Site $\rightarrow$ Room $\rightarrow$ Rack $\rightarrow$ Server $\rightarrow$ Workload), quản lý cảnh báo, sự cố và phiếu bảo trì (ticket).
    \item \textbf{Thiết kế Mô hình CSDL chuẩn hóa:} Xây dựng sơ đồ thực thể liên kết (ERD), chuyển sang mô hình quan hệ và thực hiện chuẩn hóa dữ liệu đạt dạng chuẩn 3NF nhằm triệt tiêu dư thừa.
    \item \textbf{Cài đặt SQL \& Truy vấn dữ liệu:} Viết câu lệnh DDL khởi tạo CSDL, DML thao tác dữ liệu và ứng dụng Đại số quan hệ để tối ưu các truy vấn báo cáo.
\end{itemize}

\textbf{Phạm vi đề tài:} Thực hiện thiết kế CSDL ở các mức Quan niệm, Logic và Vật lý phục vụ mô hình thử nghiệm Mini Data Center trong khuôn khổ bài tập lớn môn học.

\vspace{0.4em}
Mặc dù đã cố gắng áp dụng kiến thức môn học để hoàn thiện báo cáo, song do hạn chế về kinh nghiệm thực tế, bài làm khó tránh khỏi những thiếu sót. Chúng em rất mong nhận được góp ý quý báu từ Thầy/Cô để đề tài được hoàn thiện hơn.

\vspace{0.4em}
\begin{raggedleft}
\textit{Chúng em xin chân thành cảm ơn!}\\
\end{raggedleft}